\section{Discussion and Conclusion}
\label{sec:discussion}

The core reason domain-specific codecs outperform general-purpose tools is that
generic compressors operate on byte streams and therefore treat the two bytes of
each 16-bit pixel as unrelated, breaking the spatial correlation that spans pixel
boundaries in big-endian encoding. Our codecs operate on 16-bit values directly,
predict each pixel from its decoded neighbours, and entropy-code only the residual,
yielding a reduction of $0.87$--$1.12$\,bpp over gzip-6 depending on the codec.

The predictor turns out to be the single most important design choice. Mode
selection on the eight benchmark images reveals that global mean subtraction
dominates on bias frames and flat fields, where a near-constant DC offset accounts
for most of the signal. The three-neighbour average prevails on smooth stellar
backgrounds, while least-squares modes are only selected for images with bright
point sources and extended spatial gradients. The bias pedestal is therefore the
most compressible feature of astronomical frames, and least-squares predictors add
value only when residual spatial structure remains non-trivial after DC removal.

The hi/lo byte split is equally important. After a good prediction, the high byte
is near-uniformly zero, yielding a near-zero-entropy stream that is often
compressed to a 2-byte single-symbol FSE packet, while all useful variation is
carried by the low byte. Conditioning the low-byte model on whether the high byte
is zero exploits this bimodal structure. speed-focused achieves the same effect
through eight adaptive context models without transmitting a frequency table,
trading a modest $0.11$\,bpp in compression ratio for a $2.3\times$ speed
advantage over all-rounder.

The results also show that compression-focused's exhaustive search over predictors,
block sizes, and context strategies recovers only an additional $0.03$\,bpp over
all-rounder, at the cost of a much longer compression pass. This confirms that
most of the remaining gap to the information-theoretic limit lies in predictor
quality rather than entropy coding, and that the fast adaptive strategy of
all-rounder already captures the majority of the available gain.

There are limitations worth noting. The corpus covers only 8 images from a single
instrument, and rankings may shift on sensors with different noise characteristics
or gain settings. Runtime measurements reflect a single wall-clock run and include
process-launch overhead for the external tools. Despite these caveats, the
consistent margin across all eight images gives confidence that the gains are
structural rather than incidental. Future work could explore residual
autocorrelation removal and 16-bit FSE tables for high-correlation blocks, where
a small additional gain appears tractable.
